\documentclass[10pt,a4paper]{amsart}
\usepackage[margin=19mm]{geometry}
\usepackage{amsmath,amssymb,mathtools}
\usepackage[scr=rsfs,cal=euler]{mathalfa}
\usepackage{xfrac}
\usepackage[unicode,hidelinks,bookmarks=false]{hyperref}
\newcommand{\ie}{i.e.\ }
\newcommand{\cf}{cf.\ }
\newcommand{\resp}{respectively\ }
\newcommand{\Subsection}{\subsection}
\providecommand{\nobreakdash}{\nobreak\hbox{-}}
\setlength{\emergencystretch}{2em}
\multlinegap0pt
\usepackage{mathtools}
\usepackage[normalem]{ulem}
\usepackage{float}
\usepackage{xcolor}
\usepackage{xfrac}
\usepackage{amssymb}
\usepackage{bm}
\newtheorem{lem}{Lemma}
\newtheorem{fact}{Fact}
\title{Symmetry-enriched topological order and quasifractonic behavior in $\mathbb{Z}_N$ stabilizer codes}
\author{Siyu He, Hao Song}
\date{Source: arXiv:2511.04430v2 (CC BY 4.0)}
\begin{document}
\maketitle

\noindent\emph{Source and licence.} Symmetry-enriched topological order and quasifractonic behavior in $\mathbb{Z}_N$ stabilizer codes. Version arXiv:2511.04430v2 (CC BY 4.0), distributed by arXiv under the Creative Commons Attribution 4.0 International licence, http://creativecommons.org/licenses/by/4.0/, as recorded in the arXiv licence metadata for this exact version and shown on the abstract page https://arxiv.org/abs/2511.04430v2. Full licence notice: \texttt{licenses/arxiv-2511.04430-CC-BY-4.0.txt}.\par\smallskip

\noindent\textit{Editorial scope.} Counted unit: the article's lem lem:local\_operator\_can\_be\_created (source0.tex lines 891--893). The article's complete original proof of it is delivered with it (source0.tex lines 895--941, one \texttt{proof} environment of 47 lines). No mathematics is written, repaired or re-derived: the statement and the proof are the article's own lines, in the article's own notation, and no retained line is substituted or re-flowed.  Retained uncounted as the source-local context the proof needs: lines 945--947.  Nothing was omitted from any retained span, so the package carries no omission record.  No retained line contains a citation, so the wrapper carries no bibliography and no bibliography entry.  No retained line contains a figure environment, an image-inclusion command or drawing code, so no image is delivered and the package declares no asset dependency.  Editorial formatting only, in addition: an AMS \texttt{amsart} preamble that restates the article's own declarations and macros used by the retained lines, namely the environments \texttt{lem}, \texttt{fact} on separate counters, together with the standard packages those lines need.  The retained environments print in this document as Lemma 1, Fact 1 in order of appearance, whereas the article prints the same items with the numbering of its own full document.  The complete native source of the article as distributed in the arXiv e-print for this exact version is bundled unchanged as \texttt{sources/188588/source0.tex}.
\begin{fact} \label{fact1}
For any non-zero Laurent polynomial $f$ in $d$ variables, a weight vector $\mathbf{w}=(w_1, \dots, w_d)$ with positive integer components can be chosen to define a weighted degree under which $f$ has a unique term of minimal degree.
\end{fact}

\begin{lem}
    For a $\mathbf{BB}_N(f,g)$ code, the vector $(1,0)^T$ corresponds to a legitimate excitation if and only if $p\nmid f$ or $p\nmid g$.\label{lem:local_operator_can_be_created}
\end{lem}

\begin{proof}
($\implies$) Assume that $\left(1,0\right)^{T}$ represents a legitimate
excitation. By definition, this means there are formal Laurent series $\alpha$ and $\beta$ (in variables $x$ and $y$) such that
\begin{equation}
1=\alpha f+\beta g.
\end{equation}
(A finite series, i.e., a polynomial, corresponds to a local operator, while an infinite series corresponds to an infinitely extended operator.) 

We proceed by contradiction. Assume that the conclusion is false, i.e.,
that $p\mid f$ and $p\mid g$. Then $p$  also divides $\alpha f+\beta g$.
This implies that $p$ divides $1$. This is a contradiction, as $p$ is not a unit in $\mathbb{Z}_{p^k}$.

($\impliedby$) Without loss of generality, assume $p\nmid f$. We
will show that $\left(1,0\right)^{T}$ is a legitimate excitation
by constructing a formal Laurent series $\alpha$ such that $\alpha f=1$.

First, we can use the division algorithm to write $f=pq+r$, where
$q$ and $r$ are Laurent polynomials. Since $p\nmid f$, the remainder
$r$ is non-zero. Since $p^{k}=0$ in the ring $\mathbb{Z}_{p^{k}}$,
expanding $r^{k}=\left(f-pq\right)^{k}$ shows that $r^{k}$ is divisible
by $f$. This means there exists a Laurent polynomial $\xi\in\mathbb{Z}_{p^{k}}\left[x^{\pm},y^{\pm}\right]$
such that 
\begin{equation}
r^{k}=f\cdot\xi.\label{eq:rk}
\end{equation}

Next, according to Fact~\ref{fact1} (given below), it is possible to choose a weighted
degree such that the lowest-degree term in $r$ is unique; let us
call it $cx^{a_{1}}y^{a_{2}}$. The lowest-degree term of $r^{k}$
is then $c^{k}x^{ka_{1}}y^{ka_{2}}$. By construction, $p\nmid c$,
so $c$ is invertible in $\mathbb{Z}_{p^{k}}$. 
Dividing both sides of Eq.~\eqref{eq:rk} by $c^{k}x^{ka_{1}}y^{ka_{2}}$ gives
\begin{equation}
1-u=f\cdot h, \label{eq:fh}
\end{equation}
where $u$ and $h$ denote polynomials defined as
\begin{equation}
u\coloneqq1-c^{-k}x^{-ka_{1}}y^{-ka_{2}}\cdot r^{k},\qquad h\coloneqq c^{-k}x^{-ka_{1}}y^{-ka_{2}}\xi.
\end{equation}
By construction, each term in $u$ has a strictly positive weighted
degree. Thus, 
$
\alpha\coloneqq h\sum_{i=0}^{\infty}u^{i}    
$
is a well-defined formal Laurent series, and Eq.~\eqref{eq:fh} implies $\alpha f=1$. Taking $\beta=0$, we have $\alpha f+\beta g=1$, which proves that $\left(1,0\right)^{T}$
is a legitimate excitation.
\end{proof}

\end{document}
